\documentclass[11pt, letterpaper]{article}

% Packages:
\usepackage[
    ignoreheadfoot, % set margins without considering header and footer
    top=2 cm, % seperation between body and page edge from the top
    bottom=2 cm, % seperation between body and page edge from the bottom
    left=2 cm, % seperation between body and page edge from the left
    right=2 cm, % seperation between body and page edge from the right
    footskip=1.0 cm, % seperation between body and footer
% showframe % for debugging
]{geometry} % for adjusting page geometry
\usepackage{titlesec} % for customizing section titles
\usepackage{tabularx} % for making tables with fixed width columns
\usepackage{array} % tabularx requires this
\usepackage[dvipsnames]{xcolor} % for coloring text
\definecolor{primaryColor}{RGB}{0, 0, 0} % define primary color
\usepackage{enumitem} % for customizing lists
\usepackage{fontawesome5} % for using icons
\usepackage{amsmath} % for math
\usepackage[
    pdftitle={Pierre-Stephane Voltaire's CV},
    pdfauthor={Pierre-Stephane Voltaire},
    pdfcreator={LaTeX with RenderCV},
    colorlinks=true,
    urlcolor=primaryColor
]{hyperref} % for links, metadata and bookmarks
\usepackage[pscoord]{eso-pic} % for floating text on the page
\usepackage{calc} % for calculating lengths
\usepackage{bookmark} % for bookmarks
\usepackage{lastpage} % for getting the total number of pages
\usepackage{changepage} % for one column entries (adjustwidth environment)
\usepackage{paracol} % for two and three column entries
\usepackage{ifthen} % for conditional statements
\usepackage{needspace} % for avoiding page brake right after the section title
\usepackage{iftex} % check if engine is pdflatex, xetex or luatex

% Ensure that generate pdf is machine readable/ATS parsable:
\ifPDFTeX
\input{glyphtounicode}
\pdfgentounicode=1
\usepackage[T1]{fontenc}
\usepackage[utf8]{inputenc}
\usepackage{lmodern}
\fi

\usepackage{charter}

% Some settings:
\raggedright
\AtBeginEnvironment{adjustwidth}{\partopsep0pt} % remove space before adjustwidth environment
\pagestyle{empty} % no header or footer
\setcounter{secnumdepth}{0} % no section numbering
\setlength{\parindent}{0pt} % no indentation
\setlength{\topskip}{0pt} % no top skip
\setlength{\columnsep}{0.15cm} % set column seperation
\pagenumbering{gobble} % no page numbering

\titleformat{\section}{\needspace{4\baselineskip}\bfseries\fontsize{12pt}{12pt}\selectfont}{}{0pt}{}[\vspace{1pt}\titlerule]

\titlespacing{\section}{
% left space:
    -1pt
}{
% top space:
    0.3 cm
}{
% bottom space:
    0.2 cm
} % section title spacing

\renewcommand\labelitemi{$\vcenter{\hbox{\small$\bullet$}}$} % custom bullet points
\newenvironment{highlights}{
\begin{itemize}[
topsep=0.10 cm,
parsep=0.10 cm,
partopsep=0pt,
itemsep=0pt,
leftmargin=0 cm + 20pt
]
}{
\end{itemize}
} % new environment for highlights


\newenvironment{highlightsforbulletentries}{
\begin{itemize}[
topsep=0.10 cm,
parsep=0.10 cm,
partopsep=0pt,
itemsep=0pt,
leftmargin=10pt
]
}{
\end{itemize}
} % new environment for highlights for bullet entries

\newenvironment{onecolentry}{
\begin{adjustwidth}{
0 cm + 0.00001 cm
}{
0 cm + 0.00001 cm
}
}{
\end{adjustwidth}
} % new environment for one column entries

\newenvironment{twocolentry}[2][]{
\onecolentry
\def\secondColumn{#2}
\setcolumnwidth{\fill, 4.5 cm}
\begin{paracol}{2}
}{
\switchcolumn \raggedleft \secondColumn
\end{paracol}
\endonecolentry
} % new environment for two column entries

\newenvironment{threecolentry}[3][]{
\onecolentry
\def\thirdColumn{#3}
\setcolumnwidth{, \fill, 4.5 cm}
\begin{paracol}{3}
{\raggedright #2} \switchcolumn
}{
\switchcolumn \raggedleft \thirdColumn
\end{paracol}
\endonecolentry
} % new environment for three column entries

\newenvironment{header}{
\setlength{\topsep}{0pt}\par\kern\topsep\centering\linespread{1.5}
}{
\par\kern\topsep
} % new environment for the header

\newcommand{\placelastupdatedtext}{% \placetextbox{<horizontal pos>}{<vertical pos>}{<stuff>}
\AddToShipoutPictureFG*{% Add <stuff> to current page foreground
\put(
\LenToUnit{\paperwidth-2 cm-0 cm+0.05cm},
\LenToUnit{\paperheight-1.0 cm}
){\vtop{{\null}\makebox[0pt][c]{
\small\color{gray}\textit{Last updated in September 2024}\hspace{\widthof{Last updated in September 2024}}
}}}%
}%
}%

% save the original href command in a new command:
\let\hrefWithoutArrow\href

% new command for external links:


\begin{document}
\newcommand{\AND}{\unskip
\cleaders\copy\ANDbox\hskip\wd\ANDbox
\ignorespaces
}
\newsavebox\ANDbox
\sbox\ANDbox{$|$}

\begin{header}
\fontsize{15 pt}{15 pt}\selectfont Pierre-Stéphane Voltaire

\vspace{2 pt}

\normalsize
\kern 5.0 pt%
\mbox{\hrefWithoutArrow{mailto:psvoltaire96@gmail.com}{psvoltaire96@gmail.com}}%
\kern 5.0 pt%
\AND%
\kern 5.0 pt%
\mbox{\hrefWithoutArrow{tel:+1-905-616-0236}{+1(905)-616-0236}}%
\kern 5.0 pt%
\AND%
\kern 5.0 pt%
\mbox{\hrefWithoutArrow{https://psvoltaire.ca/}{psvoltaire.ca}}%
\kern 5.0 pt%
\AND%
\kern 5.0 pt%
\mbox{\hrefWithoutArrow{https://www.linkedin.com/in/pierre-stephane-voltaire/}{linkedin.com/in/pierre-stephane-voltaire}}%
\kern 5.0 pt%
\AND%
\kern 5.0 pt%
\mbox{\hrefWithoutArrow{https://github.com/PierreStephaneVoltaire}{github.com/PierreStephaneVoltaire}}%
\end{header}

\vspace{2 pt - 0.3 cm}

\section{Summary}

\begin{onecolentry}
\begin{highlightsforbulletentries}

Senior DevOps Engineer with over 7 years of experience architecting cloud-native infrastructure, implementing Infrastructure as Code, and building scalable CI/CD pipelines. Strong coding background with proficiency in \textbf{Python}, \textbf{Go}, \textbf{JavaScript}, and \textbf{Java} combined with extensive IaC expertise using \textbf{Terraform}, \textbf{AWS CDK}, and \textbf{CloudFormation}. Has hands-on experience across \textbf{AWS} and expanding expertise in \textbf{Azure}.
% , including serverless technologies (\textbf{Lambda}, \textbf{Step Functions}), containerization (\textbf{Docker}, \textbf{Kubernetes}), and high-availability architectures.
% Adept at collaborating with development teams to identify and resolve application bottlenecks, facilitating cloud-friendly refactoring initiatives.

\end{highlightsforbulletentries}
\end{onecolentry}

\section{Skills}

\begin{onecolentry}
\textbf{Programming \& Scripting:} Java,
TypeScript, JavaScript, Golang, Python, Bash, Powershell, SQL
\end{onecolentry}

\vspace{0.05 cm}
\begin{onecolentry}
\textbf{Development library \& Frameworks:} Spring Boot, Flask, Gin, Vue.js, Angular, ReactJS
\end{onecolentry}

\vspace{0.05 cm}
\begin{onecolentry}
\textbf{CI\//CD:} CodeDeploy, CodePipeline, CodeBuild, CircleCI, GitHub Actions, Jenkins, GitLab CI, ArgoCD, Hashicorp Vault \end{onecolentry}

\vspace{0.05 cm}
\begin{onecolentry}
\textbf{Infrastructure as code:} Terraform, Terragrunt, AWS CDK, Cloudformation, Crossplane \end{onecolentry}

\vspace{0.05 cm}
\begin{onecolentry}
\textbf{Containerization \& Orchestration:} Docker, Podman, Docker Compose, Helm,
ECS, Fargate, Kubernetes
\end{onecolentry}



\section{Experience}

\begin{twocolentry}{
Oct 2025 – Present
}
\textbf{Senior DevOps Specialist}
\end{twocolentry}

\begin{onecolentry}
\textit{Nasdaq} | \textit{Toronto, ON, CA}
\end{onecolentry}

\vspace{0.10 cm}
\begin{onecolentry}
\begin{highlights}
\item \textbf{Maintained and enhanced a cloud-native PHP/Python platform} running on \textbf{AWS Lambda}, diagnosing and resolving production issues to ensure platform reliability and continuity of service.
\vspace{0.05 cm}

\item \textbf{Owned end-to-end Terraform infrastructure stacks} provisioning AWS environments, with Ansible, Packer, and Python as core automation layers; participated in incident response for AWS performance and networking failures.
\vspace{0.05 cm}

\item \textbf{Built and maintained Jenkins CI/CD pipelines} for platform delivery and initiated an AI productivity program using GitHub Copilot instruction files to standardize AI-assisted development workflows and documentation.
\vspace{0.05 cm}

\item \textbf{Designed event-driven scheduling capabilities} using Amazon S3, EventBridge, and SQS, enabling asynchronous processing for an enterprise regulatory reporting and risk management SaaS platform.
\vspace{0.05 cm}

\item \textbf{Implemented bidirectional data synchronization} between Amazon S3 and Azure Blob Storage, enabling reliable cross-cloud data exchange between AWS and Azure environments.
\end{highlights}
\end{onecolentry}

\vspace{0.2 cm}

\begin{twocolentry}{
Mar 2021 – Oct 2025
}
\textbf{Senior Platform Engineer}
\end{twocolentry}

\begin{onecolentry}
\textit{Mitsubishi UFJ Financial Group Investor Service} | \textit{Toronto, ON, CA}
\end{onecolentry}

\vspace{0.10 cm}
\begin{onecolentry}
\begin{highlights}
\item  \textbf{Built AI-powered merge request reviewer} using Python, Tree-sitter and AWS Bedrock, \textbf{reducing manual review burden and improving code quality standards}.

\item \textbf{Developed cloud-native applications} including a FastAPI (Python) GitLab backup solution protecting \textbf{$\sim$1,000 repositories} against extended outages, with deployment orchestrated through ArgoCD and Kustomize for flexible environment management. 
\vspace{0.05 cm}

\item \textbf{Automated EKS cluster upgrades} for \textbf{20+ production clusters} using GitLab CI and Terraform, eliminating manual operations and \textbf{reducing AWS extended support costs by \$6K/month}.  
\vspace{0.05 cm}

\item \textbf{Architected and implemented CI/CD pipelines} using GitLab CI, CodeDeploy, and custom Python/Bash/Powershell automation scripts, \textbf{reducing production deployment time by 94\%} from hours to minutes.  
\vspace{0.05 cm}

\item \textbf{Developed} secure and efficient \textbf{processes and tools for managing secrets and injecting application} configurations using Jenkins and Lambda functions in Python and Node.js.
\vspace{0.05 cm}

\item \textbf{Led security hardening initiatives} during the migration of Java and .NET web applications to AWS, \textbf{aligning with SOC 2 compliance} by enhancing network security with HTTPS and ACM certificates, encrypting EFS storage with AWS KMS, and utilizing IAM for access control enforcement.
\vspace{0.05 cm}

\item \textbf{Migrated legacy source control systems} from TFS and SVN to GitLab, ensuring seamless version control transition and implementing GitOps workflows. 
\vspace{0.05 cm}

\item  \textbf{Optimized artifact management} across Sonatype Nexus, GitLab, and Amazon ECR for Java, .NET, Python packages and Docker images, improving deployment reliability.  \vspace{0.05 cm}

\item \textbf{Provisioned cloud environments} for greenfield applications using Infrastructure as Code tools including CloudFormation, Terraform, and AWS CDK, enabling rapid scalable deployments.
\vspace{0.05 cm}

\item \textbf{Enhanced observability} by integrating Datadog APM into native cloud applications, \textbf{ accelerating troubleshooting} and improving performance visibility. 
\vspace{0.05 cm}
\end{highlights}
\end{onecolentry}

\vspace{0.2 cm}

\begin{twocolentry}{
Aug 2020 – Mar 2021
}
\textbf{FullStack Developer}
\end{twocolentry}

\begin{onecolentry}
\textit{Airborne App} | \textit{Oakville, ON, CA}
\end{onecolentry}

\vspace{0.10 cm}
\begin{onecolentry}
\begin{highlights}


\item  \textbf{Developed microservices} in Go that integrate with multiple AWS services including S3, Lambda, RDS, and API Gateway, \textbf{focusing on scalable backend architecture}. 
\vspace{0.05 cm}

\item \textbf{Built automated CI/CD pipelines} leveraging GitHub Actions and Terraform to automatically lint code, execute unit and integration tests for pull requests, and \textbf{reduce production bug introduction}. 

\item  Created ephemeral testing environments using GitHub Actions, Terraform, and S3, \textbf{accelerating frontend validation processes} and streamlining collaboration between developers and product owners. 
\end{highlights}
\end{onecolentry}

\vspace{0.2 cm}

\begin{twocolentry}{
May 2019 – Apr 2020
}
\textbf{FullStack Developer}
\end{twocolentry}

\begin{onecolentry}
\textit{EllisDon} | \textit{Mississauga, ON, CA}
\end{onecolentry}

\vspace{0.10 cm}
\begin{onecolentry}
\begin{highlights}

\item \textbf{Developed and maintained full-stack applications} using ReactJS, Spring Boot (Java), Express (Node.js), NestJS (TypeScript), and MongoDB, ensuring seamless feature integration and enhanced code maintainability.
\vspace{0.05 cm}

\item \textbf{Implemented automated CI pipeline}s for unit testing and code quality enforcement using linting tools, improving project maintainability and reducing technical debt.
\vspace{0.05 cm}
\end{highlights}
\end{onecolentry}

\section{Certifications} 

\begin{twocolentry}{
June 2025 – June 2027

}
\hrefWithoutArrow{https://www.credly.com/badges/c2c2db13-a312-48b3-b4b2-e720b3b13c7b/linked_in_profile}{\textbf{HashiCorp Certified: Terraform Associate (003) (HCTA0-003)}}
\end{twocolentry}
\vspace{0.10 cm}

\begin{twocolentry}{
In Progress
}
\textbf {AWS Certified Solutions Architect - Professional (SAP-C02)}\end{twocolentry}

\vspace{0.10 cm}

\section{Education}

\begin{twocolentry}{Sept 2016 – Apr 2020}
\textbf{Sheridan College}, Bachelor of Computer Science – Mobile Computing \end{twocolentry}

\vspace{0.10 cm}

\begin{twocolentry}{Sept 2014 – Aug 2016}
\textbf{Sheridan College}, Computer Systems Technician – Software Engineering\end{twocolentry}

\vspace{0.10 cm}

\end{document}
